\chapter{The AutoML Pipeline}
\label{ch:pipeline}

Chapter~\ref{ch:arch} explained \emph{how} components communicate;
this chapter explains \emph{what} the platform actually does with a
dataset. The pipeline is written as if the LLM issues each tool call
manually, because in practice the LLM composes them in exactly this
order.

\section{Pipeline Overview}
\label{sec:pipeline-overview}

\begin{figure}[H]
\centering
\begin{tikzpicture}[node distance=0.5cm and 0.5cm, font=\small\sffamily]
  \node[al-node, minimum width=2.6cm]                              (raw){Raw CSV / \\ TSV / Parquet};
  \node[al-node, right=of raw, minimum width=2.6cm]                (ing){ingest \\ \& profile};
  \node[al-node, right=of ing, minimum width=2.6cm]                (val){pandera \\ validation};
  \node[al-node, right=of val, minimum width=2.6cm]                (eda){full EDA \\ + plots};
  \node[al-node, below=1cm of eda, minimum width=2.6cm]            (bak){parallel \\ bake-off};
  \node[al-node, left=of bak, minimum width=2.6cm]                 (hpo){Optuna \\ sweep (opt.)};
  \node[al-node, left=of hpo, minimum width=2.6cm]                 (shap){SHAP or \\ feature imp.};
  \node[al-node, left=of shap, minimum width=2.6cm]                (exp){notebook \\ + PDF};

  \draw[al-flow] (raw)--(ing);
  \draw[al-flow] (ing)--(val);
  \draw[al-flow] (val)--(eda);
  \draw[al-flow] (eda.south) |- (bak.east);
  \draw[al-flow] (bak)--(hpo);
  \draw[al-flow] (hpo)--(shap);
  \draw[al-flow] (shap)--(exp);
\end{tikzpicture}
\caption{End-to-end tabular AutoML pipeline realised by \projname.}
\label{fig:pipeline}
\end{figure}

Each box maps to one or two MCP tools; the arrows are optional to the
extent that the LLM may skip stages (e.g.\ if the user only asks for
EDA). \impl{} labels apply to every stage --- there is no ``proposed
but unimplemented'' step in Figure~\ref{fig:pipeline}.

\section{Stage 1: Ingestion}
\label{sec:stage-ingest}

The user uploads a file via the Streamlit sidebar; bytes are handed
to \code{workspace.save\_uploaded\_dataset}, which writes to
\file{data/uploads/<safe\_name>}.

The LLM then invokes \tool{ingest\_dataset(file\_path)}. The tool
delegates to \code{workspace.resolve\_dataset} to obtain the local
path, reads the file with \code{pandas.read\_csv}
(or \code{read\_parquet} / \code{read\_json}), and returns:

\begin{itemize}[leftmargin=1.4em]
  \item File name, format, and shape.
  \item Memory footprint in KB.
  \item Column names and dtypes.
  \item Head preview (default 5 rows).
  \item Total missing count and duplicate-row count.
\end{itemize}

This response is small enough to fit comfortably in the LLM context
and gives the model everything it needs to reason about subsequent
questions --- ``does the target column exist?'', ``is this
classification or regression?'', ``what fraction is missing?''.

\section{Stage 2: Schema Validation}
\label{sec:stage-validate}

\tool{validate\_schema\_with\_pandera(file\_path, expected\_columns,
strict\_dtypes)} infers a pandera schema and validates the frame
against it in lazy mode. Rules are constructed per column as follows:

\begin{itemize}[leftmargin=1.4em]
  \item Numeric columns receive
        \code{pa.Check.in\_range(min, max)}.
  \item Low-cardinality string columns ($\le 20$ uniques) receive
        \code{pa.Check.isin(allowed\_categories)}.
  \item Nullability follows from observed missing values.
\end{itemize}

The tool returns the inferred rules, whether validation succeeded,
and up to 50 failure cases. This is used by the LLM to warn the user
about outliers, unexpected string values, or off-range numeric
readings before any modelling is done.

\section{Stage 3: Exploratory Data Analysis}
\label{sec:stage-eda}

\tool{run\_full\_eda(file\_path, target\_column)} produces:

\begin{enumerate}[leftmargin=1.4em]
  \item Descriptive statistics per numeric column: mean, std, min,
        Q25, median, Q75, max, IQR, skewness, IQR-outlier count.
  \item Descriptive statistics per categorical column: unique count,
        top value, top frequency, top-5 values.
  \item Missing-value summary per column and overall.
  \item Three matplotlib plots saved as PNG artefacts:
        \begin{itemize}[leftmargin=1.2em]
          \item Correlation matrix (\emph{if} $\ge 2$ numeric cols).
          \item Small multiple of histograms for the first
                12 numeric columns.
          \item Bar chart of missing percentages (\emph{if} any
                missing values exist).
        \end{itemize}
  \item Optional target analysis if the user supplied a
        \code{target\_column}.
\end{enumerate}

\tool{render\_correlation\_matrix} produces a standalone correlation
heatmap with Pearson, Spearman, or Kendall as requested, and reports
the top pairs with $|\rho| \ge 0.6$.

\section{Stage 4: Model Selection --- The Bake-Off}
\label{sec:stage-bakeoff}

The bake-off is the analytical heart of the platform.
\tool{run\_parallel\_bake\_off(file\_path, target\_column, test\_size,
cv, random\_state)} performs the following procedure.

\begin{algorithm}[H]
\caption{Parallel bake-off (\code{mcp\_servers/modeling.py}).}
\label{alg:bakeoff}
\begin{algorithmic}[1]
\Require dataset $\mathcal{D}$, target $t$, test size $s$, folds $k$, seed $\sigma$
\State drop rows with NaN in $t$; abort if $t \notin \mathcal{D}$
\State $y \gets \mathcal{D}[t]$; if $y$ non-numeric $\Rightarrow$ label-encode
\State detect problem type from $y$ (classification if $|\text{unique}| \le \max(20, 0.05n)$ or non-numeric; else regression)
\State build \code{ColumnTransformer}: median-impute+scale numerics, most-frequent+one-hot categoricals
\State stratified split $(X_{tr}, X_{te}, y_{tr}, y_{te})$
\State $\mathcal{M} \gets \{\text{RandomForest}, \text{XGBoost}, \text{LightGBM}, \text{baseline}\}$ (LR or Ridge)
\ForAll{$m \in \mathcal{M}$ in parallel (thread pool, max 4)}
  \State $\text{cv}_m \gets \operatorname{cross\_val\_score}(\text{Pipe}(pre, m), X_{tr}, y_{tr}, cv=k, \text{scoring}=\text{primary})$
  \State refit on $(X_{tr}, y_{tr})$; predict on $X_{te}$
  \State compute accuracy \& $F_1$ (classification) or RMSE \& $R^2$ (regression); record train time
\EndFor
\State sort leaderboard by primary metric descending; abort if all failed
\State $\text{champion} \gets \arg\max$
\State persist champion pipeline (joblib) and $X_{tr}[0{:}500]$ sample (parquet) as artefacts
\State \Return leaderboard, champion, model\_id, x\_train\_sample\_id
\end{algorithmic}
\end{algorithm}

The primary metric is F1-weighted for classification and $R^2$ for
regression. The choice of a small candidate set (four models) is
deliberate: it favours interpretability and short training times
over exhaustive search, which is appropriate for the ``upload a CSV
and get a baseline'' use case.

\section{Stage 5: Hyperparameter Optimisation (Optional)}
\label{sec:stage-hpo}

\tool{trigger\_hyperparameter\_sweep(file\_path, target\_column,
model\_family, time\_budget\_seconds, cv, random\_state)} performs a
TPE-based Optuna study on one of RandomForest, XGBoost, or LightGBM.
The search space is defined inside the objective function
(``define-by-run''):

\begin{itemize}[leftmargin=1.4em]
  \item \textbf{RandomForest:} \code{n\_estimators} $\in [100, 600]$,
        \code{max\_depth} $\in [3, 30]$, \code{min\_samples\_split}
        $\in [2, 20]$.
  \item \textbf{XGBoost:} \code{n\_estimators} $\in [100, 600]$,
        \code{max\_depth} $\in [3, 12]$, \code{learning\_rate}
        $\in [10^{-2}, 3\!\cdot\!10^{-1}]$ (log-uniform),
        \code{subsample} and \code{colsample\_bytree} $\in [0.6, 1.0]$.
  \item \textbf{LightGBM:} \code{n\_estimators} $\in [100, 600]$,
        \code{num\_leaves} $\in [15, 127]$, \code{learning\_rate}
        $\in [10^{-2}, 3\!\cdot\!10^{-1}]$ (log-uniform),
        \code{feature\_fraction} $\in [0.6, 1.0]$.
\end{itemize}

The study is bounded by wall-clock time; progress is streamed at
2-second intervals with the best value observed so far.

\section{Stage 6: Explainability}
\label{sec:stage-explain}

The champion model is loaded by ID via
\code{workspace.artifact\_path}. \tool{calculate\_shap\_values}
picks an explainer based on the class name of the underlying
estimator:

\begin{itemize}[leftmargin=1.4em]
  \item Tree-family (\code{Forest}, \code{XGB}, \code{LGBM},
        \code{GradientBoosting}) $\rightarrow$ \code{TreeExplainer}
        \citep{lundberg2020treeshap}.
  \item Otherwise $\rightarrow$ \code{KernelExplainer} against a
        100-sample background.
\end{itemize}

The tool renders \code{shap.summary\_plot} and stores it as a plot
artefact, and returns a top-20 mean-absolute-SHAP ranking.

\tool{generate\_feature\_importance\_plot} is a lighter alternative
that uses the estimator's native \code{feature\_importances\_}
attribute (available on tree-based models) and produces a horizontal
bar chart of the top-$k$ features.

\section{Stage 7: Export --- Reproducible Notebook and PDF}
\label{sec:stage-export}

\tool{generate\_jupyter\_notebook} builds a notebook dictionary in
memory following the nbformat 4 schema. The generated cells cover:

\begin{enumerate}[leftmargin=1.4em]
  \item Header with project metadata.
  \item Imports (adjusted to the champion family).
  \item Data load.
  \item Train/test split.
  \item Preprocessing pipeline construction.
  \item Champion model fit.
  \item Evaluation with the appropriate metrics.
  \item \code{joblib.dump} of the champion.
\end{enumerate}

The notebook is written to disk with \code{json.dumps} and persisted
as a \code{training\_notebook} artefact. A user can download it,
run it in Jupyter, and reproduce the pipeline offline.

\tool{compile\_pdf\_report} uses ReportLab \citep{reportlab} to
produce a Letter-format PDF containing a title block, dataset and
target metadata, the leaderboard table, and optional notes.

\section{Putting the Pipeline Together}
\label{sec:pipeline-full}

A typical single-turn request from a user
--- ``run a full EDA on my data, train a model to predict
\code{<target>}, explain it, and generate a notebook and PDF'' ---
results in the tool sequence shown in
Table~\ref{tab:tool-sequence}.

\begin{table}[H]
\centering
\caption{Typical tool sequence issued by the LLM in response to a
compound request.}
\label{tab:tool-sequence}
\begin{tabular}{@{}rll@{}}
\toprule
\# & \textbf{Tool} & \textbf{Purpose} \\
\midrule
1 & \tool{list\_uploaded\_files}         & Confirm active dataset exists \\
2 & \tool{ingest\_dataset}               & Shape/dtypes preview \\
3 & \tool{run\_full\_eda}                & Statistics + 3 plots \\
4 & \tool{render\_correlation\_matrix}   & Additional heatmap \\
5 & \tool{run\_parallel\_bake\_off}      & Champion + leaderboard \\
6 & \tool{calculate\_shap\_values}       & Global feature importances \\
7 & \tool{generate\_jupyter\_notebook}   & Reproducible notebook \\
8 & \tool{compile\_pdf\_report}          & PDF summary \\
\bottomrule
\end{tabular}
\end{table}

This sequence has been observed in the actual application during
smoke-testing of the migrated codebase; the corresponding server logs
were retained during development to confirm each artefact was
generated.
